\section*{Anexos}
\addcontentsline{toc}{section}{Anexos}

\subsection*{Anexo A. Registro de evidencias y cadena de custodia}
\addcontentsline{toc}{subsection}{Anexo A. Registro de evidencias y cadena de custodia}

Conserve la evidencia original sin modificaciones y analice una copia de trabajo. Registre el hash SHA-256 de la copia original antes de abrirla con cualquier herramienta de análisis.

\begin{landscape}
	\begin{table}[H]
		\caption{Registro de evidencias y cadena de custodia}
		\label{tab:evidencias}
		\centering
		\footnotesize
		\renewcommand{\arraystretch}{1.35}
		\setlength{\tabcolsep}{4pt}
		\begin{tabularx}{\textwidth}{@{}>{\raggedright\arraybackslash}p{0.15\textwidth} >{\raggedright\arraybackslash}p{0.09\textwidth} >{\raggedright\arraybackslash}p{0.08\textwidth} >{\raggedright\arraybackslash}p{0.11\textwidth} >{\ttfamily\scriptsize\raggedright\arraybackslash}p{0.32\textwidth} >{\raggedright\arraybackslash}p{0.10\textwidth}@{}}
			\toprule
			\textbf{Archivo} & \textbf{Tipo} & \textbf{Tamaño} & \textbf{Fecha/hora} & \textbf{SHA-256} & \textbf{Analista} \\
			\midrule
			& & & & & \\ \hline
			& & & & & \\ \hline
			& & & & & \\ \hline
			& & & & & \\ \hline
			& & & & & \\ \hline
			& & & & & \\ \hline
			& & & & & \\ \hline
			& & & & & \\
			\bottomrule
		\end{tabularx}
	\end{table}
\end{landscape}

\subsection*{Anexo B. Lista de verificación de entrega}
\addcontentsline{toc}{subsection}{Anexo B. Lista de verificación de entrega}

\begin{table}[H]
	\caption{Lista de verificación de entrega}
	\label{tab:checklist-entrega}
	\centering
	\footnotesize
	\renewcommand{\arraystretch}{1.3}
	\begin{tabularx}{\textwidth}{@{}>{\raggedright\arraybackslash}p{0.36\textwidth} c Y >{\raggedright\arraybackslash}p{0.14\textwidth}@{}}
		\toprule
		\textbf{Entregable} & \textbf{¿Incluido?} & \textbf{Observaciones} & \textbf{Responsable} \\
		\midrule
		Informe técnico con resumen ejecutivo y alcance. & \checkbox & & \\
		Diez fichas de caso completas. & \checkbox & & \\
		Capturas legibles y PCAP de laboratorio. & \checkbox & & \\
		Hashes SHA-256 y cadena de custodia. & \checkbox & & \\
		Reglas de detección propuestas. & \checkbox & & \\
		Conclusiones y limitaciones. & \checkbox & & \\
		Código fuente y README de la Parte~II. & \checkbox & & \\
		Diagrama de arquitectura y modelo de datos. & \checkbox & & \\
		Demostración con datos sintéticos. & \checkbox & & \\
		Política de privacidad y consentimiento. & \checkbox & & \\
		\bottomrule
	\end{tabularx}
\end{table}

\subsection*{Anexo C. Glosario}
\addcontentsline{toc}{subsection}{Anexo C. Glosario}

\begin{table}[H]
	\caption{Términos frecuentes del taller}
	\label{tab:glosario}
	\centering
	\footnotesize
	\renewcommand{\arraystretch}{1.2}
	\begin{tabularx}{\textwidth}{@{}>{\raggedright\arraybackslash}p{0.22\textwidth} Y@{}}
		\toprule
		\textbf{Término} & \textbf{Definición} \\
		\midrule
		Anomalía & Desviación observable respecto de una línea base o comportamiento esperado. \\
		E2EE & Cifrado de extremo a extremo: solo los participantes pueden leer el contenido. \\
		Falso positivo & Alerta que señala un comportamiento legítimo como si fuera malicioso. \\
		Geofencing & Definición de una zona geográfica y generación de alertas al entrar o salir de ella. \\
		Hash SHA-256 & Huella digital de 256 bits que permite verificar la integridad de un archivo. \\
		IDS / IPS & Sistema de detección / prevención de intrusiones. \\
		IOC & Indicador de compromiso: señal observable que sugiere actividad maliciosa. \\
		MITM & Ataque de intermediario que intercepta o modifica una comunicación. \\
		PCAP & Archivo de captura de paquetes de red. \\
		SIEM & Plataforma de correlación y gestión de eventos e incidentes de seguridad. \\
		TLS & Protocolo de seguridad para el transporte de datos en red. \\
		& \\
		\bottomrule
	\end{tabularx}
\end{table}
